
\chapter[Numerics in Meteorology]{Numerics in Meteorology: Past, Present and Future}

\textbf{Author:} Hilary Weller

\section*{Abstract}

Meteorology is the science of weather and climate. The value of meteorology
is prediction, and the most powerful predictions use large models
that numerically solve the equations governing the dynamics and physics
of the atmosphere and ocean. At the heart of these models are dynamical
cores. The dynamical core solves the partial differential equations
that predict winds, temperature, pressure and the moisture content
of the atmosphere. This chapter will review some of the numerical
methods used in the dynamical core of models of the atmosphere, focusing
on the requirements and constraints, and the advantages and disadvantages
of different methods and the different grids of the sphere.

\section{The Early Days of Numerical Weather Prediction}

\subsection{L. F. Richardson's First Forecast}

The first attempt at numerical weather prediction was the heroic,
complex, but naive hand calculations by \citet{Rich22}, who used
finite differences to solve what we now call the three dimensional
primitive equations. I say naive because with hindsight we can see
that his methods would fail, even with modern super-computing. But
Richardson did not have the tools of scale analysis and numerical
analysis that we have now.

The primitive equations give the rate of change of the wind vector,
$\mathbf{u}$, the air density, $\rho$, and the potential temperature,
$\theta$:
\begin{align}
\frac{D\mathbf{u}}{Dt}+2\bm{\Omega}\times\mathbf{u}+\mathbf{g} & =-\frac{1}{\rho}\nabla p\label{eq:mom}\\
\frac{D\rho}{Dt} & =-\rho\nabla\cdot\mathbf{u}\label{eq:cont}\\
\frac{D\theta}{Dt} & =0.\label{eq:theta}
\end{align}
where $\bm{\Omega}$ is the rotation rate of the planet and $\mathbf{g}$
is the acceleration due to gravity. The total derivatives, or advection
operators of a quantity $\phi$ are defined as
\begin{equation}
\frac{D\phi}{Dt}=\frac{\partial\phi}{\partial t}+\mathbf{u}\cdot\nabla\phi\label{eq:linAdvect}
\end{equation}
which represents the rate of change of $\phi$ following with the
wind, rather than the rate of change at a fixed point, $\partial\phi/\partial t$.
Advection the movement by the wind. The potential temperature is a
wonderful meteorological concept
\begin{equation}
\theta=T\left(p_{0}/p\right)^{R/c_{p}}
\end{equation}
where $T$ is the temperature, $p_{0}=1000\text{ hPa}$, $R$ is the
specific gas constant of dry air and $c_{p}$ is the specific heat
capacity at constant pressure. It is the temperature that an air parcel
would reach (meteorologists love air parcels), if it were brought
adiabatically to 1000~hPa. That is, a small volume of air in the
upper atmosphere is brought down to near the ground, while contracting
as the pressure increases, without any other exchanges of energy between
the air parcel and its surroundings, so it heats up as it descends,
but its potential temperature, $\theta$, remains the same. One of
the wonderful aspects of potential temperature is that its transport
equation (\ref{eq:theta}), unlike that of temperature, does not contain
a pressure work term, that describes how temperature changes when
the air expands or contracts, which can make the equations more difficult
to solve. Another joy of potential temperature is that it informative
about the stability of the atmosphere. Temperature always decreases
with height in the troposphere (below the stratosphere). But potential
temperature usually increases or is constant. If potential temperature
is ever found to decrease with height then the atmosphere will be
unstable, updraughts will form, which can lead to clouds and precipitation.

There is an additional unknown in the primitive equation set (\ref{eq:mom}-\ref{eq:theta}),
the pressure, $p$. Unlike the dependent variables ($\mathbf{u}$,
$\rho$ and $\theta$), this cannot evolve independently. It is tied
to density and temperature by an equation of state, such as the ideal
gas law
\begin{equation}
p=\rho RT.
\end{equation}

The equations that Richardson solved were more complex; he included
an equation for total water content, he included turbulent (eddy)
viscosity, he wrote the equations in component form in spherical coordinates
(in terms of latitude and longitude), he included a source term for
the $\theta$ equation due to solar radiation and longwave cooling,
he included complex surface processes, including things like evaporation
from foliage, and he made the excellent hydrostatic assumption
\begin{equation}
\frac{\partial p}{\partial z}=-\rho g
\end{equation}
where $z$ is the local vertical direction and $g$ is the magnitude
of the acceleration due to gravity. This can be used, with the continuity
equation (\ref{eq:cont}) to eliminate the vertical velocity, $w$. 

Richardson was optimistic because he classified weather prediction
as a ``marching'' problem (or hyperbolic), with information propagating
from the initial conditions and from the boundaries, as opposed to
a ``jury'' problem, in which ``jurymen'' sitting at the boundary
must all agree simultaneously on then solution (elliptic). This classification
was correct, but the information propagates so quickly, in the form
of gravity waves, that we need to use numerical methods for solving
elliptic problems (including simultaneous solutions in different locations)
for solving the primitive equations.

\citet{Rich22} used centered finite differences, knowing them to
give second-order accuracy (error are proportional to the square of
the grid spacing, so as resolution increases, errors drop quadratically).
He used a leapfrog method in time, which he described as ``step-over''.
To illustrate this method, we can consider the solution of the $\theta$
equation in just one spatial dimension, $x$, which is the one-dimensional
linear-advection equation:
\begin{equation}
\frac{\partial\theta}{\partial t}+u\frac{\partial\theta}{\partial x}=0.\label{eq:linAdvect1}
\end{equation}
If we store values of $\theta$ on a grid in an array, indexed by
$j$, then $\theta_{j}$ is $\theta$ at position $j\Delta x$ where
$\Delta x$ is the grid spacing. Then we can make the second-order,
centered approximation
\begin{equation}
\frac{\partial\theta}{\partial x}\bigg|_{j}=\frac{\theta_{j+1}-\theta_{j-1}}{2\Delta x}.\label{eq:CS}
\end{equation}
We also define $\theta^{(n)}$ to be $\theta$ at time level $n$
which is at time $n\Delta t$, where $\Delta t$ is the time step.
We can make a similar centered approximation for the rate of change:
\begin{equation}
\frac{\partial\theta}{\partial t}\bigg|^{(n)}=\frac{\theta^{(n+1)}-\theta^{(n-1)}}{2\Delta t}.\label{eq:CT}
\end{equation}
Equations (\ref{eq:CS}) and (\ref{eq:CT}) can be substituted into
(\ref{eq:linAdvect1}) to give an equation for the next time level
based on old values:
\begin{equation}
\theta_{j}^{(n+1)}=\theta_{j}^{(n-1)}-\frac{u\Delta t}{\Delta x}\left(\theta_{j+1}^{(n)}-\theta_{j-1}^{(n)}\right).\label{eq:CTCSadvect}
\end{equation}
This is the explicit, centered in time, centered in space (CTCS) advection
scheme (explicit meaning that the solution at the next time level
is found from existing solutions). \citet{Rich22} described a few
alternatives for what to do at the very first time step, when you
do not know values from two previous time levels, only one. This cost
effective scheme is not used much as it generates solutions with oscillations.
This is worth trying ...

\subsection{Exercise: Numerical Solution of the Linear Advection Equation}\label{subsec:Ex_CTCSadvection}

Solve the linear advection equation (\ref{eq:linAdvect1}) using the
centered in time, centered in space scheme (CTCS, \ref{eq:CTCSadvect}).
Use a periodic domain $x\in[0,1)$ of 40 points. Test the code using
the initial conditions
\begin{equation}
\theta\left(x,0\right)=\begin{cases}
\frac{1}{2}\left(1-\cos\frac{2\pi\left(x-0.1\right)}{0.3}\right) & \text{for }x\in\left[0.1,0.4\right]\\
1 & \text{for }x\in\left[0.7,0.9\right]\\
0 & \text{otherwise}
\end{cases}
\end{equation}
The Courant number is
\begin{equation}
c=\frac{u\Delta t}{\Delta x}.\label{eq:Courant}
\end{equation}
Set the Courant number to be 0.4 and run for 10 time steps. You can
compare with the analytic solution
\begin{equation}
\theta\left(x,t\right)=\theta\left(x-ut,0\right).
\end{equation}
Can you animate the solution? What can you say about the part of the
solution that started from the smooth initial conditions? What do
the discontinuities do to the solution? Why?

Next try setting the Courant number to be 1.1 and animate the solution
for 5 time steps. What is happening now? Why?

\subsection{Insights Enabling Numerical Weather Prediction}

Jule \citet{Char48} explained why Richardson's forecast produced
such wild results and, in doing so, was able to find the right simplification
of the equations to make them tractable by numerical solution, with
smooth solutions. Charney pointed out that aerodynamics had become
tractable, in part by assuming that the air is close to incompressible,
so the continuity equation (\ref{eq:cont}) is simplified:
\begin{equation}
\nabla\cdot\mathbf{u}=0\implies\frac{D\rho}{Dt}=0.\label{eq:ico}
\end{equation}
This simplification is not useful for weather prediction because density
varies by orders of magnitude between the surface and the stratosphere.
The primitive equations (\ref{eq:mom}-\ref{eq:theta}) include the
representation of acoustic (sound) waves and gravity waves. Deep gravity
waves, in which large regions of the atmosphere move up and down under
the influence of gravity (like waves on the surface of the ocean)
can propagate horizontally at similar speed to acoustic waves. If
a numerical model for equations that represent these waves is explicit
(solutions for the next time level are found from previous solutions)
and if gradients are approximated using nearby grid point values,
then time steps must be small to avoid unstable oscillations. To explore
this, attempt the exercises in section \ref{subsec:solveSW}.

\subsubsection{Example: Numerical Solution of Gravity Waves}\label{subsec:solveSW}

The primitive equations can be simplified and averaged in the vertical
to produce the shallow-water equations:
\begin{align}
\frac{D\mathbf{u}}{Dt}+2\bm{\Omega}\times\mathbf{u} & =-g\nabla h\label{eq:SWu_vec}\\
\frac{Dh}{Dt} & =-h\nabla\cdot\mathbf{u}\label{eq:SWh_vec}
\end{align}
where the velocity, $\mathbf{u}$, is only non-zero in the locally
horizontal direction, ie perpendicular to $\mathbf{g}$. These equations
work in this Cartesian form on the surface of the earth, with $g=|\mathbf{g}|$,
with $\mathbf{g}$ perpendicular to the local surface whereas $\bm{\Omega}=\mathbf{k}\Omega$
where $\mathbf{k}$ points up through the North pole. If we consider
a local tangent plane to the sphere, then the shallow water equations
can be written
\begin{align}
\frac{Du}{Dt}-fv & =-g\frac{\partial h}{\partial x}\label{eq:SWu}\\
\frac{Dv}{Dt}+fu & =-g\frac{\partial h}{\partial y}\label{eq:SWv}\\
\frac{Dh}{Dt} & =-h\left(\frac{\partial u}{\partial y}+\frac{\partial v}{\partial y}\right)\label{eq:SWh}
\end{align}
where $\left(u,v\right)$ are the velocity components in the $\left(x,y\right)$
directions and the Coriolis parameter, $f=2\Omega\sin\phi$ where
$\phi$ is the latitude. In order to represent only gravity waves,
these equations are linearised by replacing $h$ with $H+h$ where
$H$ is uniform in space and time and $u,v,h$ are assumed small.
Just considering one spatial dimension and ignoring rotation ($f=0$)
gives a one-dimensional wave equation
\begin{align}
\frac{\partial u}{\partial t} & =-g\frac{\partial h}{\partial x}\label{eq:linSWu}\\
\frac{\partial h}{\partial t} & =-H\frac{\partial u}{\partial x}\label{eq:linSWh}
\end{align}
which has gravity wave speed $\sqrt{gH}$. If the mean fluid depth,
$H$, is, for example, the height of the troposphere, around 10~km,
and $g\sim10\text{ m/s}$ then the gravity wave speed is 316~m/s,
close to the speed of sound which is around 343~m/s. The linearised
shallow water equations can be solved explicitly using finite differences
on a staggered grid, with $u$ and $h$ stored at alternating locations
\citep[an excellent choice made by][]{Rich22}
\begin{align}
\frac{h_{j}^{(n+1)}-h_{j}^{(n)}}{\Delta t} & =-H\frac{u_{j+1/2}^{(n)}-u_{j-1/2}^{(n)}}{\Delta x}\label{eq:linSWh_FB_C}\\
\frac{u_{j+1/2}^{(n+1)}-u_{j+1/2}^{(n)}}{\Delta t} & =-g\frac{h_{j+1}^{(n+1)}-h_{j}^{(n+1)}}{\Delta x}\label{eq:linSWu_FB_C}
\end{align}
where $\Delta t$ is the time step and $n$ is the time level number
so that $t=n\Delta t$ is the time. The values of $u$ and $h$ are
stored in arrays indexed by an integer $j$. The array for $h$ at
index $j$ represents $h$ at location $j\Delta x$, where $\Delta x$
is the grid spacing. However the array for $u$ at grid point $j$
will actually represent $u$ at location $\left(j+1/2\right)\Delta x$.
Eqn (\ref{eq:linSWh_FB_C}) can be rearranged to give $h_{j}^{(n+1)}$
at every grid point using a forward difference in time (the time derivative
at time $n$ is approximated using time $n$ and the forward time,
$n+1$). Once we know $h_{j}^{(n+1)}$, this is used to update each
$u_{j+1/2}^{(n+1)}$ using an explicit backward time difference. The
resulting explicit time stepping scheme is called forward-backwards
\citep{Mes77} and it is stable provided the Courant number associated
with gravity waves, $c=\sqrt{gH}\Delta t/\Delta x\le1$. \citet{Rich22}
used $\Delta t=6\text{ hours}$ and $\Delta x=200\text{ km}$. 

\subsubsection*{Exercises}
\begin{enumerate}
\item Work out the gravity wave Courant number of Richardson's simulation,
assuming a height of 10~km. 
\item Work out what time step Richardson would have needed to treat gravity
waves stably.
\item Write code to solve the one-dimensional linearised shallow-water equations
using the forward-backward time-stepping on a staggered grid (\ref{eq:linSWh_FB_C},\ref{eq:linSWu_FB_C})
in the periodic domain $x\in[0,L)$ using the setup:
\[
\begin{array}{c}
\Delta x=2\times10^{5}\text{ m},\ g=10\text{ m/s},\ L=8\times10^{6}\text{ m},\ H=10\times10^{3}\text{ m}\\
\begin{array}{cc}
u\left(x,t=0\right),\  & h\left(x,t=0\right)=\begin{cases}
\frac{1}{2}\left(1-\cos\frac{6\pi\left(x-L/3\right)}{L}\right) & \text{for }x\in\left[L/3,\ 2L/3\right]\\
0 & \text{otherwise}
\end{cases}\end{array}
\end{array}
\]
Animate the results for 5 time steps using Richardson's time step
of 6~hours. Compare with a stable time step running for the same
total duration (so you will need many more, much smaller time steps).
\item The analytic solution returns to the initial conditions after one
complete revolution of the periodic domain. Run your model for exactly
one revolution by choosing a time step so that the number of time
steps is twice the number of points in space. 
\begin{enumerate}
\item Once you have a simulation that returns approximately to the initial
conditions, you can calculate the normalised root mean square error
of $h$ (the $\ell_{2}$ error norm):
\begin{equation}
\ell_{2}=\sqrt{\sum_{j=0}^{n_{x}-1}\frac{\left(h_{j}-h\left(x_{j}\right)\right)^{2}}{\left(h\left(x_{j}\right)\right)^{2}}}
\end{equation}
where $h_{j}$ is the numerical solution and $h\left(x_{j}\right)$
is the analytic solution at $x_{j}$. For this setup, you should get
at $\ell_{2}$ error norm of around 0.04.
\item Double the number of grid points and time steps (so that the Courant
number stays the same). Repeat the simulation, with the solution returning
to the initial conditions. Recalculate the $\ell_{2}$ error norm.
This is a second-order method so $\ell_{2}$ should reduce by a factor
of 4.
\end{enumerate}
\item Write code to solve the shallow-water equations a co-located method
(without grid-staggering), with $h_{j}$ and $u_{j}$ stored at the
same locations:
\begin{align}
\frac{h_{j}^{(n+1)}-h_{j}^{(n)}}{\Delta t} & =-H\frac{u_{j+1}^{(n)}-u_{j-1}^{(n)}}{2\Delta x}\label{eq:linSWh_FB_A}\\
\frac{u_{j}^{(n+1)}-u_{j}^{(n)}}{\Delta t} & =-g\frac{h_{j+1}^{(n+1)}-h_{j-1}^{(n+1)}}{2\Delta x}\label{eq:linSWu_FB_A}
\end{align}

Run this code with the same parameters as the staggered grid code.
You should find:
\begin{enumerate}
\item The errors are a factor of 4 larger than when using a staggered grid.
\item You can use a time step twice as big as the staggered grid and still
get stable results.
\item If you initialise the height with a grid-scale oscillation:
\begin{equation}
h\left(x_{j},t=0\right)=\begin{cases}
1 & \text{for }j\text{ even}\\
-1 & \text{otherwise}
\end{cases}
\end{equation}
you should find a spurious stable solution. This is a computational
mode -- a solution with artefacts that do not go away as the resolution
is increased.
\item Re-run the staggered grid code with the initial conditions consisting
of a grid-scale oscillations. Does this solution seem to more accurate?
What solution would you expect?
\end{enumerate}
\end{enumerate}

\subsubsection{Solving Elliptic Equations}\label{subsec:ellipticSolutions}

The primitive equations (\ref{eq:mom}-\ref{eq:theta}), the linear
advection equations (\ref{eq:linAdvect1}) and the shallow-water equations
(\ref{eq:SWu_vec}-\ref{eq:SWh_vec}) are hyperbolic, meaning that
information propagates at a finite speed. This encouraged Richardson
in developing an explicit time-stepping method. However, the propagation
speeds are so fast that we actually need a numerical method similar
to a numerical method for solving an elliptic equations. By making
the incompressible approximation (\ref{eq:ico}) for the primitive
equations, acoustic waves have infinite speed, making the problem
elliptic. The incompressible approximation for the shallow-water equations
can also be used to derive the barotropic vorticity equation. This
is an equation for the vertical component of the total vorticity,
$\eta$, which includes the vertical component of the relative vorticity,
$\xi=\left(\nabla\times\mathbf{u}\right)\cdot\mathbf{k}$, and planetary
rotation (the Coriolis parameter, $f$):
\begin{equation}
\eta=\xi+f\label{eq:pv}
\end{equation}
Taking the curl of the momentum, applying some vector-calculus identities
and assuming incompressiblity gives the barotropic voriticity equation,
\begin{equation}
\frac{D\eta}{Dt}=0\label{eq:baroVortEqn}
\end{equation}
(see the mathematical exercises below). This seems wonderfully simple!
The difficult bit comes is calculating the velocity from the vorticity
by inverting (\ref{eq:pv}) which is exactly the kind of ``jurymen''
problem that Richardson wanted to avoid. \citet{Char48} realised
that some such elliptic solution was essential for avoiding severe
time-step restrictions. \citet{Char49} knew that weather forecasters
relied mostly on charts of pressure to predict the weather. It should
not be necessary to solve the full primitive equations, whose initialisation
relies on accurate knowledge of winds, density and temperature. The
wind is close to geostrophic balance
\begin{equation}
\begin{array}{cccc}
u\approx-\frac{1}{\rho f}\frac{\partial p}{\partial y},\ \  & v\approx\frac{1}{\rho f}\frac{\partial p}{\partial x},\ \  & w\approx0 & \text{ or }\mathbf{u}=\mathbf{k}\times\frac{\nabla p}{\rho f}\end{array}\label{eq:geoBalance}
\end{equation}
so wind can be determined from the pressure (flowing anti-clockwise,
or cyclonically, around low pressure in the Northern hemisphere).
The winds advect (move) the pressure, making it evolve slowly, making
the weather predictable. However geostrophic balance cannot be used
to predict the weather. Geostrophic balance (\ref{eq:geoBalance})
implies incompressibility and, substituting (\ref{eq:geoBalance})
into the momentum equation (\ref{eq:mom}) results in the wind and
other variables never changing. This is a persistence forecast - the
weather tomorrow will be the same as today. We don't need equations
to make a persistence forecast. So it seems that we need need some
way to ignore the ``entirely irrelevant gravity motions''. \citet{CFvN50}
made the first numerical weather predictions by finding a pressure
level in the atmosphere, $p^{*}$, where horizontal divergence is
smallest. The equations are then formulated on this pressure surface,
which has geopotential height, $z$. \citet{CFvN50} used $p^{*}=500\text{ hPa}$
and, to this day, much weather data is reported on the $500\text{ hPa}$
surface and ECMWF use the height of the $500\text{ hPa}$ surface,
$z_{500}$, to verify their forecasts. Geostrophically balanced winds
can be written in terms of $z$
\begin{equation}
\mathbf{v}=-\frac{g}{f}\mathbf{k}\times\nabla z\label{eq:geoBalanceZ}
\end{equation}
where $\mathbf{v}$ is the horizontal component of the wind. Hence
the two equations to be solved, for total vorticity, $\eta$, and
geopotential height, $z$, for the simplest weather prediction are
\begin{align}
\frac{D\eta}{Dt} & =0, & g\nabla\cdot\frac{1}{f}\nabla z & =\xi.\label{eq:baroVortGeoVort}
\end{align}


\subsubsection*{Mathematical Exercises}
\begin{enumerate}
\item From the barotropic vorticity equation (\ref{eq:baroVortEqn}), derive
the equation for the relative vorticity, $\xi$ on a $\beta$ plane
\begin{equation}
\frac{D\xi}{Dt}+\beta v=0
\end{equation}
where we assume that Coriolis varies linearly in the $y$ direction
\begin{equation}
f=f_{0}+\beta y
\end{equation}
where $f_{0}$ is a constant.
\item Eqn (\ref{eq:geoBalance}) gives geostrophic balance in terms of pressure,
$p$, and (\ref{eq:geoBalanceZ}) gives geostrophic balance in terms
of geostrophic height, $z$. Use hydrostatic balance and the chain
rule on a surface, $S$, where pressure is uniform
\begin{align*}
\frac{\partial p}{\partial x}\bigg|_{S} & =\frac{\partial p}{\partial x}\bigg|_{z}+\frac{\partial p}{\partial z}\bigg|_{x}\frac{\partial z}{\partial x}\bigg|_{S}\left(=0\right) & \frac{\partial p}{\partial z} & =-\rho g
\end{align*}
 to derive (\ref{eq:geoBalanceZ}) from (\ref{eq:geoBalance}).
\item Use geostrophic balance to derive the second equation in (\ref{eq:baroVortGeoVort})
\begin{equation}
g\nabla\cdot\frac{1}{f}\nabla z=\xi.
\end{equation}
\item Derive the barotropic vorticity equation (\ref{eq:baroVortEqn}) for
two-dimensional flow, $\mathbf{u}=\left(u,v,0\right)^{T}$, by taking
the curl of (\ref{eq:SWu_vec}), assuming incompressibilty ($\nabla\cdot\mathbf{u}=\frac{\partial u}{\partial y}+\frac{\partial v}{\partial y}=0$),
uniform $h$ and using the vector calculus identities:
\begin{itemize}
\item $\mathbf{u}\cdot\nabla\mathbf{u}=\left(\nabla\times\mathbf{u}\right)\times\mathbf{u}+\nabla\frac{|\mathbf{u}|^{2}}{2}$
for any vector, $\mathbf{u}$
\item $\nabla\times\left(\nabla K\right)=0$ for any scalar, $K$
\end{itemize}
\end{enumerate}

\subsubsection*{Coding Exercise}

Invert vorticity to find geopotential height, $z$ using a finite
difference scheme in 1D. Check your answer by calculating vorticity
from $z$. Boundary conditions and initial vorticity ...

\subsubsection{Elliptic Solutions of Hyperbolic Equations -- the semi-implicit
method}

Section \ref{subsec:ellipticSolutions} described approximations of
the hyperbolic primitive equations (\ref{eq:mom}-\ref{eq:theta})
by the elliptic vorticity equation (\ref{eq:baroVortEqn}). The real
atmosphere is, of course fully compressible, with finite speed acoustic
and gravity waves and non-hydrostatic motions, so for the most accurate
weather forecasts, all of these process need to be included, but without
the time step being limited by the speed of the fastest waves. This
is where semi-implicit time stepping comes in. The full equations
are written as linearised equations that can be solved implicitly,
and the non-linear and other slow parts are solved explicitly. Rather
than describing semi-implicit time stepping for the primitive equations,
we can consider the shallow-water equations (\ref{eq:SWu_vec}-\ref{eq:SWh_vec})
and treat the fast, gravity waves implicitly and the advection and
Coriolis terms explicitly. This description will be for a time-stepping
scheme that is first-order accurate in time, but the same method can
used with schemes that are higher-order accurate in time. Discretising
in time only:
\begin{equation}
\frac{DX}{DT}=\frac{\partial X}{\partial t}+\mathbf{u}\mathbf{\cdot}\nabla X\approx\frac{X^{(n+1)}-X^{(n)}}{\Delta t}+\mathbf{u}\mathbf{\cdot}\nabla X^{(n)}
\end{equation}
for some dependent variable $X$, where the advection term, $\mathbf{u}\mathbf{\cdot}\nabla X$
is considered slow and so is treated explicitly, so the value at the
old time level is used. So (\ref{eq:SWu_vec}-\ref{eq:SWh_vec}) become
\begin{align}
\frac{\mathbf{u}^{(n+1)}-\mathbf{u}^{(n)}}{\Delta t}+\mathbf{u}^{(n)}\mathbf{\cdot}\nabla\mathbf{u}^{(n)}+2\bm{\Omega}\times\mathbf{u}^{(n)} & =-g\nabla h^{(n+1)}\label{eq:uSI}\\
\frac{h^{(n+1)}-h^{(n)}}{\Delta t}+\mathbf{u}^{(n)}\mathbf{\cdot}\nabla h^{(n)} & =-h^{(n)}\nabla\cdot\mathbf{u}^{(n+1)}\label{eq:hSI}
\end{align}
with the terms involved in gravity-wave propagation (in the linearised
equations \ref{eq:linSWu},\ref{eq:linSWh}) treated implicitly, using
the new time level. Variables $\mathbf{u}$ and $h$ both consist
of arrays of values over a grid, and $h^{(n+1)}$ is needed to update
$\mathbf{u}$ and $\mathbf{u}^{(n+1)}$ is needed to update $h$.
To proceed, we split $\mathbf{u}^{(n+1)}$ from (\ref{eq:uSI}) into
its explicit and implicit parts and substitute this into (\ref{eq:hSI}):
\begin{align}
\mathbf{u}^{(n+1)} & =\mathbf{u}^{\prime}-\Delta tg\nabla h^{(n+1)}\label{eq:backS}\\
\text{where }\mathbf{u}^{\prime} & =\mathbf{u}^{(n)}-\Delta t\left\{ \mathbf{u}^{(n)}\mathbf{\cdot}\nabla\mathbf{u}^{(n)}+2\bm{\Omega}\times\mathbf{u}^{(n)}\right\} \label{eq:uPrimeEqn}\\
\implies\frac{h^{(n+1)}-h^{(n)}}{\Delta t} & -\Delta tgh^{(n)}\nabla^{2}h^{(n+1)}=-\mathbf{u}^{(n)}\mathbf{\cdot}\nabla h^{(n)}-h^{(n)}\nabla\cdot\mathbf{u}^{\prime}.\label{eq:hEqn}
\end{align}
The first step is to calculate $\mathbf{u}^{\prime}$ explicitly from
(\ref{eq:uPrimeEqn}). Eqn (\ref{eq:hEqn}) appears to be an elliptic
equation for $h$ and is sometimes called the pressure equation or
a Helmholtz equation. All of the unknown values of $h$ can be put
into an array. The discretisation of $\nabla^{2}h$ can be represented
as a matrix multiplied by this array. Then this matrix equation can
be solved to find the array of all of the $h$ values. The final step
is the back-substitution, to calculate $\mathbf{u}^{(n+1)}$ from
$\mathbf{u}^{\prime}$ and the new values of $h^{(n+1)}$.

\subsubsection*{Exercise}

We will assume a one-dimensional staggered grid (eqns \ref{eq:linSWu},\ref{eq:linSWh}),
without rotation. So the continuous equations are
\begin{align*}
\frac{\partial u}{\partial t}+u\frac{\partial u}{\partial x} & =-g\frac{\partial h}{\partial x}\\
\frac{\partial h}{\partial t}+u\frac{\partial h}{\partial x} & =-h\frac{\partial u}{\partial x}
\end{align*}

\begin{enumerate}
\item Write the semi-implicit method, (\ref{eq:backS}-\ref{eq:hEqn}) for
this one-dimensional system.
\item We will also assume that $u>0$ everywhere, for all time. This allows
us to use a stable upwind approximation for the advection terms. So
$u^{\prime}$ can be calculated as
\begin{equation}
u_{j+1/2}^{\prime}=u_{j+1/2}^{(n)}-\Delta t\ u_{j+1/2}^{(n)}\frac{u_{j+1/2}^{(n)}-u_{j-1/2}^{(n)}}{\Delta x}.
\end{equation}
In exercise \ref{subsec:Ex_CTCSadvection} we used a centered in space
approximation for advection, which lead to oscillations in the solution.
This upwind approximation, with information for calculating $\partial u/\partial x$
coming only from upwind, leads to smooth solutions but it is less
accurate. Using the approximation $u_{j}=\frac{1}{2}\left(u_{j-1/2}+u_{j+1/2}\right)$,
write down the upwind approximation for the term $u\frac{\partial h}{\partial x}$
at location $j$.
\item Using grid-staggering as in (\ref{eq:linSWh_FB_C}), write down the
approximation for the term $h\frac{\partial u^{\prime}}{\partial x}$
at location $j$.
\end{enumerate}
Maybe just use implicit to solve the linearised eqns and compare with
previous exercise?

\section{A Golden Era of Numerical Weather Prediction}

Through the 1950s and 1960s, simulations of the atmosphere became
more realistic and the models became more and more similar to those
that are used now. \citet{Smag58} solved the the primitive equations
in three dimensions (with just two vertical levels). But the paper
mostly discusses causes of instability and no plots of numerical results
are presented, suggesting that there was little realism in the results.
\citet{Phil59} simulated the entire northern hemisphere with one
horizontal layer, with a 24 hour simulation taking 84 minutes. Like
\citet{Rich22}, \citet{Phil59} was still blaming grid-scale oscillations
on the lack of smoothness in the initial conditions rather than on
the numerical methods. \citet{Smag63} created a two-layer model between
the equator and 64\textsuperscript{o}N which he was able to run for
60 days, simulating baroclinic instability, poleward heat transport
and a Hadley circulation, with some realism. Grid-scale oscillations
were controlled using eddy viscosity to mimic the effects of turbulent
diffusion. \citet{Smag63} mentioned Arakawa's emerging work on energy
conserving methods which would prove revolutionary in atmosphere and
ocean modelling.

\subsection{Three Dimensional Solutions on Mainframe Computers}

\subsubsection{Arakawa's Numerics}\label{subsec:ArakawaGrids}

\citet{AL77} described the numerical methods used on the UCLA general
circulation model; a landmark in the development of stable and accurate
numerical methods for the atmosphere, a 92 page article that is still
being cited (3,538 citations in total, with over 100 citations every
year since 2014). The model is three dimensional with 12 vertical
levels, $72\times45$ grid points in the horizontal and a time step
of 6 minutes. \citet{AL77} is best known for the description of the
Arakawa grids - 5 arrangements of prognostic variables on a grid,
labelled A-E. They wrote the momentum equation, for the horizontal
velocity, $\mathbf{v}$, in vector-invariant form
\begin{equation}
\frac{D\mathbf{v}}{Dt}+f\mathbf{k}\times\mathbf{v}+\frac{1}{\rho}\nabla_{h}p+\nabla_{h}K=\mathbf{F}_{h}
\end{equation}
where $K=\frac{1}{2}|\mathbf{v}|^{2}$ is the kinetic energy, $\mathbf{F}_{h}$
is the horizontal component of the frictional force and $\nabla_{h}$
is the horizontal component of the gradient. Use of the vector-invariant
form enables exact energy conservation, with losses in gravitiational
potential being exactly transferred to kinetic energy, and back. The
Arakawa grids A-E are shown in figure \ref{fig:Arakawa-A-E}. The
advantages and dis-advantages of these grids can best be seen in terms
of the linearised shallow water equations. The one-dimensional linearised
shallow water equations without Coriolis are (\ref{eq:linSWu},\ref{eq:linSWh}).
The two-dimensioal linearised shallow water equations can be obtained
by replacing $D/Dt$ with $\partial/\partial t$ in (\ref{eq:SWu}-\ref{eq:SWh})
and replacing $h$ with the constant, uniform average height, $H$,
on the right hand side of (\ref{eq:SWh}). As was investigated in
example \ref{subsec:solveSW}, if all variables are located together,
as in the A-grid, then the pressure (or height) gradients needed to
update the velocity are calculated over three grid points, missing
the central location (eqn \ref{eq:linSWu_FB_A}) and the divergence
needed to update the height in (\ref{eq:SWh}) is calculated over
three points in each direction, missing the central location (eqn
\ref{eq:linSWh_FB_A}). This omission of the central location means
that grid-scale oscillations can develop -- oscillations in $h$
can be present which do not lead to height gradients, and oscillations
in $u$ and $v$ can be present which do not lead to divergence. These
oscillations are spurious computational modes and make the A-grid
problematic for atmospheric modelling. Considering the height gradients
and the divergence terms of eqns (\ref{eq:SWu}-\ref{eq:SWh}), the
C-grid appears to be ideal -- compact height gradients in the $x$-direction
(using two points without missing out a central point), are centered
at the $u$ location, compact height gradients in the $y$-direction
are at the $v$ location and the divergence can be calculated at the
$h$ location without missing any velocities at $h$. This means that
gravity wave propagation is most accurate and free of computational
modes on the C-grid, and this grid was chosen by \citet{AL77} for
the UCLA model. 

\begin{figure}
\begin{tabular}{ccccc}
A & B & C & D & E\tabularnewline
\resizebox{0.18\textwidth}{!}{\input{climate_numFigures/ArakawaGrids/A.pdftex_t}} & \resizebox{0.18\textwidth}{!}{\input{climate_numFigures/ArakawaGrids/B.pdftex_t}} & \resizebox{0.18\textwidth}{!}{\input{climate_numFigures/ArakawaGrids/C.pdftex_t}} & \resizebox{0.18\textwidth}{!}{\input{climate_numFigures/ArakawaGrids/D.pdftex_t}} & \resizebox{0.18\textwidth}{!}{\input{climate_numFigures/ArakawaGrids/E.pdftex_t}}\tabularnewline
\end{tabular}

\caption{Arakawa grids A-E on finite volumes,}\label{fig:Arakawa-A-E}
\end{figure}

Gravity waves are not the only phenomena being modelled, and even
in the linearised shallow water equations the Coriolis terms need
to be represented accurately. If we consider only the rates of change
and the Coriolis terms, we get equations for inertial oscillations
\begin{align}
\frac{\partial u}{\partial t} & =fv & \frac{\partial v}{\partial t} & =-fu\label{eq:inertialOsc}
\end{align}
which have solutions with air (or water) going round in circles with
radius $\sqrt{u_{o}^{2}+v_{0}^{2}}$ and angular frequency $f$. In
order to model (\ref{eq:inertialOsc}) most accurately, $u$ and $v$
should be at the same locations on the grid, otherwise spatial averaging
will be needed. . $u$ and $v$ are at the same location for grids
A, B and E but not C or D. The UK Met Office chose the B-grid for
their first unified model of weather and climate. For the C-grid,
with $h$ stored at grid point $\left(i,j\right)$, with $x=i\Delta x$
and $y=j\Delta y$, $u$ will be stored at $\left(i+1/2,j\right)$
and updated with averaging for the Coriolis term:
\begin{align}
u_{i+1/2,j}^{(n+1)} & =u_{i+1/2,j}^{(n)}\\
 & +\Delta t\left\{ f\frac{v_{i,j-1/2}+v_{i+1,j-1/2}+v_{i,j+1/2}+v_{i+1,j+1/2}}{4}-g\frac{h_{i+1,j}-h_{i,j}}{\Delta x}\right\} \nonumber 
\end{align}
This averaging reduces accuracy and enables spurious computational
modes at coarse resolution, but the accurate representation of gravity
waves has proven to be more important for atmospheric modelling. This
may seem counter intuitive, considering that \citet{CFvN50} described
these waves as ``entirely irrelevant gravity motions''. They are
not irrelevant. They are the mechanism by which the flow comes into
geostrophic balance -- if there is a region of high height and no
flow along the boundary of the high height, then gravity waves will
be generated, flow will be created and the fluid will move towards
geostrophic balance (a process called geostrophic adjustment). If
gravity waves are filtered then this process is assumed to happen
instantaneously. 

The D-grid may appear ideal for representing fluids that are close
to geostrophic balance, since the terms involved in geostrophic balance
are the largest in the momentum equation (for example, $v$ is in
the same place as the most compact calculation of $\partial p/\partial x$).
However these are both terms in the $u$ equation, and $u$ is at
a different location. The E-grid is the same as the B-grid but rotated
by 45\textsuperscript{o}.

\subsubsection*{Exercise}

Write a C-grid and an A-grid shallow water model?

\subsection{The Spectral Method}

The European Centre for Medium Range Weather Forecasts (ECMWF) consistently
produces the most accurate global weather forecasts, with skill improving
at the rate of one day per decade \citep{HJV+24}. ECMWF specialise
in forecasts with lead times of one to two weeks. Skillfull forecasts
at these longer lead times was made possible when they switched from
finite differences to the spectral method \citep{GJ82,SBJ+89} to
solve the primitive equations on there sphere. Spectral methods have
exponential convergence. This means that they have the highest possible
order of accuracy for the number of grid points. This comes at the
cost of all to all communication. This sounds like a big problem for
massively parallel computers. However, it turns out that ECMWF's spectral
model is continuing to scale well. 

\subsubsection{The Mathematics of the Spectral Method}

This section gives a simple example of a spectral method. \citet{Dur11}
provides a more general definition.

In order to solve
\begin{equation}
\frac{\partial\phi}{\partial t}+F\left(\phi\right)=0\label{eq:generalPDE}
\end{equation}
for some function, $F$, we represent $\phi$ as a sum of globally
defined basis functions. For simplicity, we will consider a one-dimensional
periodic domain, with $x\in[-\pi,\pi]$, so that $\phi$ can be represented
as a sum of $2N+1$ Fourier modes, $e^{ikx}$, with Fourier coefficients,
$a_{k}\left(t\right)$:
\begin{equation}
\phi\left(x,t\right)=\sum_{k=-N}^{N}a_{k}\left(t\right)e^{ikx}.\label{eq:FourierDecomp}
\end{equation}
The initial Fourier coefficients, $a_{k}\left(t_{0}\right)$, can
be found using the orthogonality property of the Fourier modes
\begin{equation}
\int_{-\pi}^{\pi}e^{ipx}e^{-iqx}dx=\begin{cases}
2\pi & \text{if }p=q\\
0 & \text{otherwise.}
\end{cases}\label{eq:FourierOrthogonal}
\end{equation}
So if (\ref{eq:FourierDecomp}) is multiplied by $e^{-iqx}$ and integrated
between $-\pi$ and $\pi$, we get equations for the initial Fourier
coefficients from the initial conditions:
\begin{align}
a_{k}\left(t_{0}\right) & =\frac{1}{2\pi}\int_{-\pi}^{\pi}\phi\left(x,t_{0}\right)e^{-ikx}dx\label{eq:FourierTransform}
\end{align}
Eqn (\ref{eq:FourierTransform}) is a Fourier transform and computers
have fast algorithms called ffts (fast Fourier transforms) for calculating
Fourier coefficients from evaluations of a function at regularly spaced
intervals. However to find solutions of (\ref{eq:generalPDE}), we
need to know how the Fourier coefficients vary in time. We need to
find each $a_{k}\left(t\right)$ so that when (\ref{eq:FourierDecomp})
is substituted into (\ref{eq:generalPDE}), (\ref{eq:generalPDE})
is satisfied as closely as possible. For the spectral method ``as
closely as possible'' means the Galerkin approximation:
\begin{equation}
\int_{-\pi}^{\pi}\left\{ \frac{\partial\phi}{\partial t}+F\left(\phi\right)\right\} e^{-ikx}dx=0\ \forall k\in-N,...,N.\label{eq:spectralGalerkin}
\end{equation}
Substituting (\ref{eq:FourierDecomp}) into (\ref{eq:spectralGalerkin}),
using the orthogonality of the Fourier modes and simplifying gives:
\begin{equation}
2\pi\frac{da_{k}}{dt}+\int_{-\pi}^{\pi}F\left(\sum_{q=-N}^{N}a_{q}\left(t\right)e^{iqx}\right)e^{-ikx}dx=0.\label{eq:spectralMethod}
\end{equation}
To make progress in evaluating the second part of the integral in
(\ref{eq:spectralMethod}), we need to know $F$. We will consider
some examples.

\subsubsection{Linear Advection using the Spectral Method}

Eqn (\ref{eq:generalPDE}) becomes the linear advection equation for
$F\left(\phi\right)=u\ \partial\phi/\partial x$ with a constant,
uniform wind speed $u$. Then (\ref{eq:spectralMethod}) simplifies
to
\begin{equation}
\frac{da_{k}}{dt}=iua_{k}\label{eq:spectralLinearAdvect}
\end{equation}
which can be solved either analytically or numerically for each $a_{k}$,
starting from the initial conditions, $a_{k}\left(t_{0}\right)$ from
the Fourier transform (\ref{eq:FourierTransform}). 

\subsubsection{Non-linear Advection using the Spectral Method}

Eqn (\ref{eq:generalPDE}) becomes Burgers equation for non-linear
advection for $F\left(\phi\right)=\frac{1}{2}\frac{\partial\phi^{2}}{\partial x}$.
Now eqn (\ref{eq:spectralMethod}) becomes
\begin{equation}
\frac{da_{k}}{dt}+\frac{ik}{2}\sum_{p+q=k}a_{p}\left(t\right)a_{q}\left(t\right)=0\label{eq:spectralBurger}
\end{equation}
which can be solved numerically using an ODE solver, starting from
the initial conditions, $a_{k}\left(t_{0}\right)$.

\subsubsection*{Exercise}

Check eqns (\ref{eq:spectralLinearAdvect}) and (\ref{eq:spectralBurger}).

Work out spectral method for linear advection with a source, sink
and diffusion.

Work out spectral method for the SWE

Write a code to solve an equation using the spectral method.

\subsection{Semi-Lagrangian Advection}\label{subsec:SL}

Until the use of semi-Lagrangian advection, time steps in weather
and climate models had been limited by the advection Courant number
(\ref{eq:Courant}) so that, typically, air or its properties cannot
be transported more than one grid box in one time step. This might
sound restrictive, but is actually useful for accuracy, so is often
not a problem. However, the state of the art until this century was
to use a latitude-longitude grid (fig \ref{fig:grids}) in which lines
of constant longitude converge towards the poles, making spatial resolution
unnecessarily high and making the time step restriction prohibitive.
Before semi-Lagrangian advection, prohibitive time step restrictions
near the poles were ameliorated using polar filtering \citep[eg.][]{CD91},
in which short wavelength waves in the zonal direction (along lines
of constant latitude) were filtered out. This enabled useful forecasts
but also removed energy from the model. The polar filtering and artificial
viscosity needed to keep the Met Office model from this era stable
were described as treacle, limiting forecast skill. This problem was
circumvented by the use of semi-Lagrangian advection \citep{Rob82,PS84,PL91,RTS+95,Hort02,DCM+05}.

The Courant-Friedrichs-Lewy (CFL) condition states that the domain
of dependence of a numerical model (the points used to calculate the
value at the next time step) must enclose the physical domain of dependence
(the location in space that actually influences the value of the next
time step if solving the equations exactly). For explicit, Eulerian
schemes, this usually means that nothing can be transported more than
one grid box in one time step. Meteorologists like to keep advection
(or transport) explicit, for least cost and best accuracy, but it
is possible, using Lagrangian approximations, to move the domain of
dependence to where it is needed. This is how semi-Lagrangian works.
Instead of approximating spatial gradients on the grid, we look upstream
to where things came from, and map the old upstream values onto the
grid at the new time level. For each grid point at the new time step,
we find a departure point, which won't fall on a grid point, and interpolate
from the old values onto the departure point. This enables long, stable
and accurate time steps, but the properties transported are not conserved.
This means that there can be spurious sources and sinks of mass, energy,
moisture, or anything else that is transported by the model. 

\begin{figure}
\resizebox{0.48\textwidth}{!}{\input{climate_numFigures/semiLagrange/semiLagrange.pdftex_t}}
\resizebox{0.48\textwidth}{!}{\input{climate_numFigures/semiLagrange/hexFlux.pdftex_t}}

\caption{Semi-Lagrangian and flux-form semi-Lagrangian}\label{fig:semiLagrangian}
\end{figure}


\subsubsection*{Exericise}

Semi-Lagrangian solution of linear advection and Burgers equation
in 1D. Look at conservation and accuracy for large time steps.

\section{Parallel Computing and the Pole Problem}

In the 1990s, numerical weather prediction used mainframe computers
such as the vector-parallel Cray YM-P. Parallel means that they had
more than one processor, running in parallel. In the 1990s, Crays
had 2, 4, 8 or of 16 processors. The atmosphere was decomposed into
a domains, one for each processor, and the equations of motion were
solved separately on each domain, with communication at the boundaries.
Each processor operated on a vector of numbers at the same time, with
vector lengths typically of 8, 16, 32 or 64. These were days of rapid
increases in computational power, with gains in efficiency due to
increases in vector length and number of processors \citep[eg.][]{BDI+95}.
With moderate parallelisation and moderate domain size, performance
was not limited by parallel communication or memory band width. These
were the day of Moore's law, with computational power doubling every
two years, which enabled steady increases in model resolution without
drastic changes to the numerical methods or algorithms.

Even in the 1990s, vector-parallel machines were being replaced by
machines with far more cheaper, off the shelf microprocessors. Even
into the 2000s and 2010s, weather prediction models scaled on these
pure parallel machines, with the atmosphere being decomposed into
more and more, smaller domains, all with overlapping halo regions
at their edges, where information is exchanged.

Around the turn of the century, a model including semi-implicit and
semi-Lagrangian methods on a C-grid was replacing the explicit B-grid
model at the UK Met Office \citep{CDM+97}, following the successful
use of SISL (semi-implicit, semi-Lagrangian) at ECMWF \citep{RTS+95}.
SISL enabled excellent weather and climate prediction on latitude-longitude
grids since ISIL overcomes time step restrictions near the poles.
At time of writing the operational model used by the UK Met Office,
EndGame, still uses SISL on a latutide-longitude grid, and has consistently
been one of the most accurate models, usually second only to ECMWF
\citep[eg.][]{HJV+24}. However the combination of further parallelisation,
increased resolution and the convergence of meridians towards the
poles means that parallelisation progress on latitude-longitude grids
is stuck -- more processors are added, domains get smaller and the
code does not run faster or even at the same speed using higher resolution.
It is not possible to overcome the large amount of communication that
must be done to create stable models with areas of stupidly high resolution
-- the weather at any location depends on the surrounding weather,
and points near the poles have a lot of grid points nearby that could
influence their weather. So inevitably there will be a lot of communication.
This communication is computationally expensive.

\section{Modern Solutions for Modern Computers}

The current ECMWF supercomputer has 7,680 AMD Epyc compute nodes (\url{https://www.ecmwf.int/en/computing/our-facilities/supercomputer-facility})
which are similar to the CPUs on desktop or laptop computers. The
National Oceanic and Atmospheric Administration (NOAA) uses a cluster
of 63,840 cores, including CPUs and GPUs (graphical processing units),
since GPUs are now being used to provide vector calculations which
are faster that the one at a time calculations on each CPU core. 

https://www.noaa.gov/organization/information-technology/hera.

The UK Met Office has recently moved to cloud based computing provided
my Microsoft's Azure service, which gives them flexibility on how
many and what type of core they use. However the weather prediction
models can rarely use the whole of any modern supercomputer because,
even if the whole computer were available, the model would not be
able to scale well enough to use so many processors efficiently; processors
would end up spending too much time waiting for data from neighbouring
processors. This is the parallel scaling problem. The current Met
Office model, on a latitude-longitude grid, can use thousands of processors
efficiently, but not tens of thousands, which is needed for the most
accurate, highest resolution solutions. In order to scale in parallel
on 10s or 100s of thousands of processor, quasi-uniform meshes of
the sphere are needed, which do not have convergence of points towards
the poles, or any other locations.

\subsection{Meshes of the Sphere}

A latitude-longitude grid (fig \ref{fig:grids}, left) is the most
obvious way to decompose the atmosphere and solve the equations of
motion using, for example, finite differences, and models using latitude-longitude
grids are producing accurate weather predictions and informative climate
projections to this day. The current Met Office global weather prediction
uses a grid of $2,560\times1,920$ points which gives grid boxes of
size $15\text{ km}\times10\text{ km}$ around the equator but at each
pole, the grid boxes will still be 10~km in the north-south direction
but less than 25~m in the east-west direction! (\url{https://www.metoffice.gov.uk/research/approach/modelling-systems/unified-model/weather-forecasting})
In order to get smooth solutions over the poles, all of the 2,560
grid points in the east-west direction around each pole will need
to be in communication, despite being on different processors. This
inevitably slows the simulation down. This century, weather prediction
has been moving away from the latitude-longitude grid due (in part)
to the parallel scaling bottle necks as the number of parallel processors
increases. Even before parallel scaling was a critical bottleneck,
ECMWF moved to a reduced grid, which coarser resolution towards the
poles \citep{HS91} in order to save in computation time and memory
requirement. This was more straightforward using a spectral model
as they key point was to use few wavenumbers close to the poles. 

\begin{figure}
\includegraphics[width=0.32\textwidth]{climate_numFigures/grids/latLon}
\includegraphics[width=0.32\textwidth]{climate_numFigures/grids/cubedSphere10}
\includegraphics[width=0.32\textwidth]{climate_numFigures/grids/hexIcos}

\caption{Grids}\label{fig:grids}
\end{figure}

\citet{ST12} reviewed meshes of the sphere, describing the challenges
in finding sufficiently accurate numerical methods for meshes such
as the cubed sphere (middle of fig \ref{fig:grids}) and icosahedral
meshes (eg. the hexagonal icosahedron, right of fig \ref{fig:grids},
which actually contains 12 pentagons). The cubed sphere in fig \ref{fig:grids}
is quasi-uniform -- all of the grid boxes have similar dimensions
and a low aspect ratio. However the mesh is not orthogonal -- local
co-ordinate systems for each of the cube faces are not at right angles.
A cubed sphere can be constructed to be orthogonal, but then resolution
clusters towards the cube corners, removing the quasi-uniform property.
Meshes of polygons (such a the hexagonal icosahedron, right of fig
\ref{fig:grids}) and meshes of triangles can be constructed to be
orthogonal. For these meshes, orthogonality needs a different definition.
The triangular icosahedron is the dual of the hexagonal icosahedron.
This means that if you put a point at the centre of each polygon of
the hexagonal icosahedron, and join the points over the edges, then
you get the triangular icosahedron. If the mesh and its dual have
edges that cross at right angles, then the mesh is orthogonal. Orthogonal
meshes have advangages in terms of accuracy and efficiency, but there
are other considerations, such as the skewness -- do the edges of
the dual grids bisect each other or cross not at the edge centres
\citep{Wel14}. 

\subsection{Numerical Methods for Alternative Meshes}

The numerical methods developed for latitude-longitude grids, such
as semi-implicit, semi-Lagrangian and the Arakawa C-grid have superb
accuracy and are efficient as long as the domain is not decomposed
onto too many parallel processors. Moving to alternative meshes, such
as the cubed sphere or icosahedra (fig. \ref{fig:grids}), means that
more complex numerical methods are needed in order to retain accuracy.
This section describes some more complex spatial discretisation.

\subsubsection{Finite volumes}\label{subsec:FV}

In their simplest form, on structured grids (such as latitude-longitude),
the finite volume method is similar to the finite difference method.
Important features of finite volume methods are:
\begin{enumerate}
\item Averages of computed quantities are stored in finite volume cells
and the averages change based on fluxes of the quantity over the surface
faces of the cell.
\item The finite volume cells can be any tessellating shape.
\item The total of the computed quantity is conserved.
\end{enumerate}
In order to solve an advection equation with the finite volume method,
the equation must first be written in flux form, with everything that
varies inside a divergence term. We can start with the continuity
equation (\ref{eq:cont}) which can be re-written in flux form using
vector-calculus identities
\begin{align}
\frac{\partial\rho}{\partial t}+\mathbf{u}\cdot\rho+\rho\nabla\cdot\mathbf{u} & =0\nonumber \\
\implies\frac{\partial\rho}{\partial t}+\nabla\cdot\left(\mathbf{u}\rho\right) & =0.\label{eq:contFlux}
\end{align}
The finite volume method predicts density weighted volume averages
(for example the mass of an atmospheric constituent, $\phi$) in each
cell. We can derive an equation for the rate of change of $\rho\phi$
using the continuity equation, the linear advection equation for $\phi$
(\ref{eq:linAdvect}) and the chain rule:
\begin{align}
\frac{\partial\rho\phi}{\partial t} & =\rho\frac{\partial\phi}{\partial t}+\phi\frac{\partial\rho}{\partial t}\nonumber \\
 & =\rho\left\{ -\mathbf{u}\cdot\nabla\phi\right\} +\phi\left\{ -\nabla\cdot\left(\mathbf{u}\rho\right)\right\} \nonumber \\
\implies\frac{\partial\rho\phi}{\partial t} & +\nabla\cdot\left(\rho\mathbf{u}\phi\right)=0.\label{eq:linAdvectFlux}
\end{align}
To simplify the definition of the finite volume method, we will assume
that $\rho$ is uniform in space and time and the atmosphere is incompressible:
$\nabla\cdot\mathbf{u}=0$. We consider a cell with volume $\mathcal{V}$
and integrate over this volume to find the rate of change of the volume
average:
\begin{equation}
\frac{\partial\tilde{\phi}}{\partial t}=\frac{1}{\mathcal{V}}\int_{\mathcal{V}}\frac{\partial\phi}{\partial t}\ dV=-\frac{1}{\mathcal{V}}\int_{\mathcal{V}}\nabla\cdot\left(\mathbf{u}\phi\right)\ dV
\end{equation}
(the last step using the linear advection equation in flux form, \ref{eq:linAdvectFlux}).
We now apply Gauss's divergence theorem to transform the integral
to an integral over the surface of the volume:
\begin{equation}
\int_{\mathcal{V}}\nabla\cdot\left(\mathbf{u}\phi\right)\ dV=\int_{S}\mathbf{u}\phi\cdot\mathbf{dS}\implies\frac{\partial\tilde{\phi}}{\partial t}=-\frac{1}{\mathcal{V}}\int_{S}\phi\mathbf{u}\cdot\mathbf{dS}
\end{equation}
where $S$ is the surface of volume $\mathcal{V}$ and $\mathbf{dS}$
is the outward pointing normal vector to the surface. If we assume
that the volume $\mathcal{V}$ has a discrete number of faces, indexed
with $f$, each with outward pointing normal vector $\mathbf{S}_{f}$
(with magnitude equal to the face area), then the finite volume discretisation
of the linear advection equation becomes
\begin{equation}
\frac{\partial\tilde{\phi}}{\partial t}=-\frac{1}{\mathcal{V}}\sum_{f}\phi_{f}\mathbf{u}_{f}\cdot\mathbf{S}_{f}\label{eq:FVadvect}
\end{equation}
where $\phi_{f}$ is the value of $\phi$ mapped onto face $f$ and
$\mathbf{u}_{f}\cdot\mathbf{S}_{f}$ is the volume flux through face
$f$, so that $\phi_{f}\mathbf{u}_{f}\cdot\mathbf{S}_{f}$ is the
volume flux of $\phi$ through face $f$. There are countless different
finite volume methods and they are distinct in the way that they map
volume averages, $\tilde{\phi}$, onto face averages, $\phi_{f}$.
In order to solve (\ref{eq:FVadvect}), the fluxes, $\mathbf{u}_{f}\cdot\mathbf{S}_{f}$,
must be supplied from elsewhere in the model. The fluxes and the face
averages, $\phi_{f}$, are equal and opposite for the two cells either
side of face $f$, so what leaves one cell will enter its neighbour,
which is what makes finite volume methods conservative, regardless
of how the mapping is done.

\subsubsection*{Exercise}

Often, on uniform and structured grids, finite volume methods can
be written as finite difference methods and vice versa. Eqn (\ref{eq:CS})
is a finite difference approximation for $\partial\phi/\partial x$
at grid point $j$. If instead we use $j$ to index one-dimensional
grid boxes, with size $\Delta x$ and edges at locations $\left(j\pm1/2\right)\Delta x$,
indexed as $j\pm1/2$, then we can derive a finite volume approximation
for $\nabla\cdot\left(\mathbf{u}\phi\right)$. If we assume that the
velocity, $\mathbf{u}=(u,0,0)$ is uniform in space and time and if
we use the centred mapping
\begin{equation}
\phi_{j+1/2}=\frac{1}{2}\left(\tilde{\phi}_{j}+\tilde{\phi}_{j+1}\right),
\end{equation}
show that the finite volume method gives the same approximation for
$\partial\phi/\partial x$ as the centred finite difference method
(\ref{eq:CS}).

\subsubsection{Flux-form Semi-Lagrangian}\label{subsec:FFSL}

The semi-Lagrangian method, described in section \ref{subsec:SL},
has played a large role in enabling unprecedented accuracy in weather
prediction. However standard (interpolating) semi-Lagrangian is not
conservative -- the total mass of the quantity advected can change
just because it is being advected. This is not correct and makes semi-Lagrangian
unsuitable for climate prediction and for the next generation in weather
prediction. A popular solution is flux-form semi-Lagrangian (FFSL),
which is a type of finite volume advection scheme. FFSL is used in
FV3 by the National Oceanic and Atmospheric Administration \citep[NOAA,][]{HCP+21},
in the NCAR Community Atmosphere Model, CAM-FV \citep{NRC+10}, in
the Japanese NICAM \citep{Miu07b}, in the German ICON model \url{https://www.dwd.de/EN/research/weatherforecasting/num_modelling/01_num_weather_prediction_modells/icon_description.html}
\citep{ZRR+15} and will be used in the next Met Office model, LFRic
\citep{BK25}. 

Rather than interpolating onto a departure point, in FFSL, interpolation
is onto the volume of air that sweeps through each face in a time
step (fig \ref{subsec:SL}, right). Then, using the finite volume
method, the contents of each cell is updated based on the amounts
that enter and leave at each face. FFSL is therefore conservative.
FFSL can be made to be stable and accurate for large time steps by
mapping onto the swept volume over multiple cells upwind of each face. 

FFSL is used so widely and is such an advantageous solution of the
linear advection problem, that it has been created independently in
different institutions and given different names. The piecewise parabolic
method \citep[PPM,][]{CW84} was one of the first FFSL methods, and
was described as a higher-order Godunov method. \citet{LLM96} described
this type of scheme as ``forward in time'', and described how to
extend them to two and three dimensions while maintaining conservation
and constancy preservation -- a field that starts constant should
remain constant when advected by a non-divergent velocity field. In
the same journal issue, \citet{LR96} described a similar scheme and
named it ``flux-form semi-Lagrangian''. This was the scheme that
is still in use in the NOAA model, FV3. There have been numerous improvements
since the 1990s, but of particular note is SWIFT \citep{BK25} which
extends the long time step stability, conservation and constancy preservation
to variable density advection, providing the advection scheme for
the new Met Office model, LFRic. 

In many ways, PPM and FFSL seem unbeatable. The only down side is
that, in order to achieve accuracy and efficiency for large time steps,
the grid must be be regular. That means that the grid must consist
of (possibly distorted) boxes that are arranged in the Cartesian directions.
This precludes the use of FFSL with large time steps for icosahedral
meshes, and means that special measures are needed to handle the corners
of the cubed-sphere mesh. This is because the COSMIC splitting that
is used to transform a scheme like PPM from one to two and three dimensions,
relies on searching for departure points along co-ordinate directions. 

\subsubsection*{Exercises}
\begin{enumerate}
\item Starting from the flux-form linear advection equation (\ref{eq:linAdvectFlux}),
prove constant preservation; that if the advected quantity, $\phi$,
starts uniform, then it will stay uniform, given the following assumptions:
\begin{enumerate}
\item The density is uniform and $\nabla\cdot\mathbf{u}=0$ (incompressible
flow).
\item The density is non-uniform and varies according to the continuity
equation (\ref{eq:contFlux}), with $\nabla\cdot\mathbf{u}\ne0$ (compressible
flow).
\end{enumerate}
\item The Lax-Wendroff advection scheme can be expressed as a FFSL scheme.
Using the same notation and one-dimensional grid as for the finite
difference CTCS scheme (eqn \ref{eq:CTCSadvect}), the Lax-Wendroff
advection scheme can be written
\begin{align*}
\frac{\phi_{j}^{(n+1)}-\phi_{j}^{(n)}}{\Delta t} & =-\frac{u_{j+1/2}\phi_{j+1/2}-u_{j-1/2}\phi_{j-1/2}}{\Delta x}\\
\text{where }\phi_{j+1/2} & =\nicefrac{1}{2}\left(1+c\right)\phi_{j}^{(n)}+\nicefrac{1}{2}\left(1-c\right)\phi_{j+1}^{(n)}\\
c & =u\Delta t/\Delta x.
\end{align*}
Write code to solve the linear advection equation using Lax-Wendroff
and compare with the CTCS solutions found in exercise \ref{subsec:Ex_CTCSadvection}.
Explore how the errors depend on Courant number for each scheme. 
\item Assuming $u$ is uniform, determine if the Lax-Wendroff scheme has
constancy preservation in one dimension. Check that your numerical
implementation also has this property.
\end{enumerate}

\subsubsection{Finite Volume A- and C-grids for unstructured meshes}

Sections \ref{subsec:FV} and \ref{subsec:FFSL} described finite
volume methods for solving the linear advection equation, which has
one dependent variable, which is stored as a volume average in each
cell. Once we have multiple variables that create waves and, in particular,
we are solving the momentum equation for the velocity components,
we must return to Arakawa's problem of where to store different variables
(section \ref{subsec:ArakawaGrids}). This has added complexity because
a finite volume mesh can include cells of arbitrary shapes, such as
the triangles, pentagons and hexagons in icosahedral meshes (fig.
\ref{fig:grids}). 

The purest form of finite volume is co-located, of A-grid, so that
the velocity is stored as volume averages in the cells, the same as
the variables that are advected by the wind. This approach is used
effectively by NICAM \citep{SMT+08} which stores all variables apart
from vertical velocity in cells. Vertical velocity is stored at the
top and bottom faces of cells. So NICAM uses the A-grid in the horizontal
and the C-grid (staggered) in the vertical. ECMWF's finite volume
model, PMAP, or FVM, currently a research model, is fully co-located
with all variable stored in cells \citep{KDK+19} and is close to
ECMWF's spectral transform model in terms of accuracy and efficiency. 

Despite excellent solutions from co-located (A-grid) finite volume
models on various grids, the A-grid is always prone to grid-scale
oscillations, which can usually be smoothed out with some kind of
diffusion or filter. This has motivated the development of C-grid
discretisations on triangles for the ICON model \citep{Bon04,ZRR+15}
and on hexagons and pentagons for, example, MPAS \citep{TRSK09,SKD+12,RPH+13}.
ICON produces skilful global weather forecasts, but the underlying
dynamical core is still prone to grid-scale oscillations and requires
diffusion. It turns out that the C-grid on triangles is not as well
behaved as hoped. \citet{Gas11} explained that the oscillations on
triangles cannot be controlled by the insufficient normal velocity
components at the triangle edges, and that the hexagonal C-grid performs
better. MPAS, the ``Model for Prediction Across Scales'', was developed
as a global version of the popular WRF model and uses the hexagonal
C-grid \citep{SKD+12}. It is free of grid-scale oscillations in height
or pressure, but spurious grid-scale circulations can appear around
the hexagon vertices, a manifestation of the Hollingsworth instability
\citep{Gas13}. 

Finite-volume methods for wave equations (such as the A- and C-grid
finite voulme methods), tend to be low-order accurate -- usually
up to second-order. However, despite the smooth, well behaved solutions,
the C-grid on the hexagonal-icosahedron is found to be much less than
second-order, with a complete loss of convergence with resolution
in some situations \citep{YWM+20}, who found serious deficiencies
in both the A- and C-grids on hexagonal icosahedral meshes. This provides
a strong motivation to move to more complex numerical methods on quasi-uniform
meshes.

\subsubsection{Spectral Elements}

\subsubsection{Discontinuous Galerkin}

\subsubsection{Mixed Finite Elements}

\section{Is AI the Future?}
